% ============================================================
% MODULO 03 - INTRODUCAO (redacao revisada; parafrase propria)
% ============================================================
\section{Introdu\c{c}\~ao}

O TDAH de apresenta\c{c}\~ao combinada re\'une desaten\c{c}\~ao e hiperatividade-impulsividade em n\'ivel mais acentuado do que o esperado para a idade, com in\'icio antes dos 12 anos, dura\c{c}\~ao de pelo menos seis meses e preju\'izo em dois ou mais ambientes de vida \cite{apa2023}. Dados de meta-regress\~ao situam a preval\^encia mundial do transtorno em torno de 5\% entre crian\c{c}as e adolescentes, ainda que os n\'umeros variem bastante conforme o m\'etodo adotado em cada estudo \cite{polanczyk2007}. Na primeira inf\^ancia, os sinais aparecem sobretudo no corpo: a crian\c{c}a corre sem parar, tem dificuldade de se envolver em brincadeiras calmas, fala em excesso, interrompe, custa a esperar a vez e perde o foco diante de qualquer est\'imulo ao redor.

Para menores de 6 anos, o manejo se apoia primeiro em recursos n\~ao medicamentosos: orienta\c{c}\~ao e treino dos pais, rotina e sono organizados, articula\c{c}\~ao entre escola e fam\'ilia e psicoterapia infantil \cite{hood2024}. A diretriz da Academia Americana de Pediatria aponta, para a faixa de 4 a 6 anos, o treinamento comportamental parental como porta de entrada, antes de qualquer medica\c{c}\~ao \cite{aap2019}. Na Abordagem Centrada na Pessoa \cite{rogers2020} e na ludoterapia dela derivada \cite{axline1984}, o trabalho cl\'inico n\~ao mira a supress\~ao do sintoma: oferece congru\^encia, aceita\c{c}\~ao positiva incondicional e compreens\~ao emp\'atica, deixando que a crian\c{c}a conduza o processo pelo brincar, dentro de limites de realidade claros.

Tr\^es lentes te\'oricas orientam a leitura do caso. A Teoria da Restaura\c{c}\~ao da Aten\c{c}\~ao prop\~oe que ambientes naturais capturam a aten\c{c}\~ao sem exigir esfor\c{c}o, dando descanso \`a aten\c{c}\~ao dirigida -- exatamente a mais sobrecarregada no TDAH \cite{hood2024,kuo2004}. O modelo bioecol\'ogico de Bronfenbrenner \cite{bronfenbrenner1996} lembra que a crian\c{c}a se desenvolve em sistemas encaixados: o microssistema do brincar, o mesossistema fam\'ilia-vizinhan\c{c}a e o exossistema da comunidade se influenciam mutuamente. A Abordagem Centrada na Pessoa, por fim, sustenta que todo organismo tende a se atualizar quando o ambiente o acolhe sem condi\c{c}\~oes.

O caso que disparou este trabalho chegou assim: M., 5 anos, com laudo de TDAH misto e hiperatividade t\~ao intensa que a menina n\~ao conseguia sequer assistir \`a televis\~ao. A fam\'ilia mudou de endere\c{c}o; no novo lugar, a crian\c{c}a passou a brincar diariamente na rua com vizinhos de idade pr\'oxima. A m\~ae relatou, e a terapeuta confirmou nas sess\~oes, que ela ficou visivelmente mais tranquila. A pergunta que da\'i nasceu \'e tamb\'em a lacuna desta revis\~ao: quase todos os estudos sobre natureza e TDAH foram feitos com escolares de 7 a 12 anos, em parques ou programas planejados, e nenhum sob a \'otica humanista brasileira. A hip\'otese de trabalho -- tratada como hip\'otese, n\~ao como fato -- \'e que a vizinhan\c{c}a segura, com pares e movimento livre, operou como um ``tempo verde-social'' di\'ario, com efeito restaurador amplificado pelo v\'inculo terap\^eutico.
